Use only the provided celestial object pointer to point at the target
specified in the questions. The marker (examiner) at the observation station
will mark the answers directly in the question sheet.

\begin{description}
\item[(1.1)] Move the pointer along the celestial equator. \hfill(1 point)
\item[(1.2)] Aim the pointer at the vernal equinox. \hfill(1 point)
\item[(1.3)] In the constellation of Pegasus and its vicinity there is an
  obvious square of bright stars (Great Square of Pegasus); aim the pointer
  at the brightest star of the square. \hfill(2 points)
\item[(1.4)] Aim the pointer at the star named alpha-Arietis
  ($\alpha$-Ari). \hfill(2 points)
\item[(1.5)] Start from the star named Aldebaran ($\alpha$-Tauri) in the
  constellation Taurus, turn the pointer 35 degrees northward followed by 6
  degrees westward (in equatorial coordinates). Then, aim the pointer at the
  brightest star in the field of view. \hfill(4 points)
\end{description}
